\section{\texttt{invokedynamic}: concatenação e lambdas}
\label{sec:invokedynamic}

As outras instruções de chamada sabem, em tempo de compilação, qual
método chamar. O \texttt{invokedynamic} adia essa decisão: na primeira
execução, a JVM chama um \emph{método de bootstrap}, que devolve o código
a usar naquele ponto; nas execuções seguintes, a mesma instrução reutiliza
o resultado, sem chamar o bootstrap de novo (JVMS §5.4.3.6).

\paragraph{Concatenação de strings (JDK~9+).}
Na listagem de \texttt{ConcatenacaoStrings}
(seção~\ref{sec:javap-basico}), \texttt{"Olá, " + nome + "! Idade: " +
idade} virou \emph{uma} instrução \texttt{invokedynamic}, de nome
\texttt{makeConcatWithConstants} e descritor
\nolinkurl{(Ljava/lang/String;I)Ljava/lang/String;}. Só os
valores variáveis vão na pilha; o texto fixo vai como receita no
argumento do bootstrap, com \texttt{\textbackslash u0001} marcando onde
entra cada valor (veja \texttt{PoolConstantes}, na seção anterior). O
\citetitle{doc:jep280} explica o motivo: até o JDK~8 o \texttt{javac}
gerava uma cadeia de \texttt{StringBuilder.append}, e as otimizações do
JIT para reconhecer essa cadeia eram frágeis. Com \texttt{invokedynamic},
a estratégia fica na biblioteca do JDK e pode melhorar sem recompilar o
seu código.

\emph{Bug ou mito:} o conselho \enquote{troque \texttt{+} por
\texttt{StringBuilder}} vem da era do JDK~8. Numa expressão única, o
\texttt{+} já é uma chamada só. Dentro de um laço que acumula numa
\texttt{String}, cada volta ainda cria uma \texttt{String} nova, e aí o
\texttt{StringBuilder} continua sendo a escolha (o tema do item 63 do
\citetitle{livro:bloch2017}).

\paragraph{Lambdas.}
\lstinputlisting[language=Java]{generated/source/LambdaInvokeDynamic.java}

\inputbytecode{LambdaInvokeDynamic}

\begin{itemize}
\item O corpo do lambda virou um método privado e sintético,
  \texttt{lambda\$dobro\$0}. Sem \texttt{-p}, o \texttt{javap} não o
  mostra.
\item Em \texttt{dobro}, uma única instrução: \texttt{invokedynamic}
  \nolinkurl{apply:()Ljava/util/function/Function;}. Com \texttt{-v}, o bootstrap
  aparece como \texttt{LambdaMetafactory.metafactory}, que recebe uma
  referência a \texttt{lambda\$dobro\$0}.
\item Não existe \texttt{LambdaInvokeDynamic\$1.class} no JAR, como
  existiria com uma classe anônima. Segundo o \citetitle{doc:jep371}
  (JDK~15), a \texttt{LambdaMetafactory} define em tempo de execução uma
  \emph{hidden class} que implementa a interface funcional.
\item Dentro de \texttt{lambda\$dobro\$0}: \texttt{intValue}, \texttt{imul},
  \texttt{Integer.valueOf}. Um \texttt{Function<Integer,\allowbreak Integer>}
  desembrulha e reembrulha a cada chamada; a seção seguinte mostra o
  custo e o risco disso.
\end{itemize}

\paragraph{Ponte com Quarkus.}
Executáveis nativos do Quarkus são gerados pelo GraalVM Native Image,
que trabalha com a hipótese de mundo fechado: todo o código é conhecido
em tempo de build. O \citetitle{doc:graalvm-compat} diz que
\texttt{invokedynamic} e method handles em geral são incompatíveis com
essa hipótese, porque podem introduzir chamadas em tempo de execução, mas
que os usos gerados pelo \texttt{javac} para lambdas e concatenação de
strings são suportados, porque não mudam o método chamado. Ou seja:
lambdas e \texttt{+} estão liberados no nativo; geração dinâmica de
código com method handles, não. A Parte~III detalha.
